% Source: https://artofproblemsolving.com/community/c3h1757873_list_of_logical_and_interesting_problems
% Stable identity: aops-topic-1757873-q05
Let $\theta$ be the angle between the hands. By the law of cosines the distance between the tips is
\[
d=\sqrt{3^2+4^2-2\cdot3\cdot4\cos\theta}=\sqrt{25-24\cos\theta}.
\]

\begin{asy}
usepackage("amsmath");
size(4.5cm);
defaultpen(fontsize(10pt));
pair O=(0,0);
draw(circle(O,4.6),linewidth(0.9));
for(int k=0;k<12;++k) draw(4.3*dir(90-30*k)--4.6*dir(90-30*k));
pair H=3*dir(70),M=4*dir(-10);
draw(O--H,linewidth(1.6));draw(O--M,linewidth(1));
draw(H--M,red+dashed);
draw(arc(O,0.9,-10,70));label("$\theta$",1.3*dir(30));
dot(O);
label("$3$",H/2,W);label("$4$",M/2,S);label("$d$",(H+M)/2,NE,red);
label("\mbox{}",(min(currentpicture,user=true)+max(currentpicture,user=true))/2,black); // keeps the text after the figure black when LaTeX typesets the labels inline
\end{asy}

As the angle from the hour hand to the minute hand goes once round, from $0^\circ$ to $360^\circ$, $d$ increases from $1$ to $7$ (at $180^\circ$) and decreases back to $1$. So in each full turn the integer values $2,3,4,5,6$ occur twice each and $1$ and $7$ once each: $12$ times.

The minute hand gains $330^\circ$ per hour on the hour hand, so in the $12$ hours it gains $12\cdot330^\circ=11\cdot360^\circ$: eleven full turns of the relative angle. At $1{:}00$ (start and end) the angle is $30^\circ$ and $d^2=25-12\sqrt3$ is not an integer, so no occurrence is double-counted at the endpoints. The answer is $11\cdot12=\boxed{132}$.
